% META: {"id":"1NSI-TC-EX-048","chapitre":"1NSI-TYPES-CONSTRUITS","type_objet":"exercice","capacites":["1NSI-TYPES-CONSTRUITS-C5"],"parcours":1,"competences":["analyser"],"duree_min":8,"mode_creation":"ex_nihilo","sources_inspiration":[],"status":"approved","origin":"needs_review","status_history":["needs_review","audited","accepted","approved"],"release_acceptance":"RELEASE_OWNER_BATCH_ACCEPTANCE","acceptance_closure_digest":"sha256:9b3ccf9a81c5520fbb7e03b7b2d3e7bf2057a02834b6d908bf3be5f49f3b553f","acceptance_receipt_digest":"sha256:4fad86f0282e0fa4e0419cca1ad6379ae7af5a280c04a4a90880f90a1c044242"}
% BEGIN-TRACE
% def ajouter(tab, val):
%     tab.append(val)
% scores = [10, 20]
% ajouter(scores, 30)
% print(scores)
% EXPECTED
% [10, 20, 30]
% END-TRACE
\begin{exercice}{1NSI-TC-EX-048}{1}{8}
On considère le programme suivant :
\begin{python}
def ajouter(tab, val):
    tab.append(val)

scores = [10, 20]
ajouter(scores, 30)
print(scores)
\end{python}
\begin{enumerate}
  \item Donner, sans exécuter, l'affichage produit.
  \item La fonction \lstinline{ajouter} ne contient pas de \lstinline{return}.
        Comment se fait-il que la liste \lstinline{scores} soit modifiée ?
  \item On parle d'\textbf{effet de bord}. Expliquer ce terme.
\end{enumerate}
\end{exercice}
